% Hướng dẫn cho Ban trình bày (BTK). Dịch: pdflatex slides.tex (hai lần), sau scripts/test-all.sh
\documentclass[aspectratio=169,12pt]{beamer}
\usepackage{../pi-hd}
\newcommand\vaitro{Hướng dẫn cho Ban trình bày}
\begin{document}

\bia{Ban trình bày}{Nhận tệp đề bài của số báo đã khoá để dàn trang.}

\chuanbi
\dungthu

\chu{Việc của Ban trình bày}{%
\begin{buoc}
\item Khi Tổng biên tập khoá kỳ, số báo có \textbf{tệp .tex chỉ có đề bài}, đã đánh số in, mức B trước, mức A sau, kèm dòng tác giả.
\item Thầy cô tải tệp ở trang \textbf{Bảng chọn bài}.
\item Hình của đề nằm trong \textbf{gói chế bản} (tệp .tex và thư mục pic/) trên Drive — Quản trị gửi đường dẫn,
  hoặc gửi gói đã dựng sẵn bản PDF xem trước.
\item Tệp dùng mẫu định dạng của Pi (\texttt{dinhdang.tex}); không có lời giải, ghi chú hay liên hệ.
\end{buoc}}

\manhinh{btk}{Tải tệp .tex}

\gopy

\chu{Khi có sự cố}{%
\begin{itemize}\setlength\itemsep{2mm}
\item \textbf{Không thấy nút Tải tệp .tex}: số báo chưa khoá — hỏi Tổng biên tập.
\item \textbf{Thiếu hình}: báo Quản trị (hình chưa có trong thư mục hình khi khoá).
\item \textbf{Chữ tiếng Việt lỗi}: tệp là UTF-8 dạng NFC; mở bằng trình soạn thảo hỗ trợ UTF-8.
\end{itemize}}
\end{document}
